\section{Conclusion}
\label{sec:conclusion}

We presented AnomalyClaw, a training-free framework for cross-domain visual anomaly detection built on adversarial debate between an Anomaly Proposer and a Normality Advocate, grounded by an extensible pool of lightweight vision experts and orchestrated by an autonomous controller.
The adversarial structure provides built-in self-verification---every anomaly claim must survive explicit counter-arguments before acceptance---addressing the hallucination and lack-of-verification problems that limit single-pass VLM approaches.

Evaluated on 1{,}440 test items across 12 diverse domains, AnomalyClaw achieves \tbd{X.XXX} macro AUROC, outperforming single-pass VLM baselines (\tbd{X.XXX}), expert-only methods (\tbd{X.XXX}), and symmetric multi-agent debate (\tbd{X.XXX}).
The framework generalizes across multiple VLM backends without modification, and the autonomous controller adapts reasoning depth to query difficulty, reducing cost on easy cases while investing more computation on ambiguous ones.

\paragraph{Limitations.}
Performance on medical domains (particularly liver CT) remains below clinical requirements, reflecting fundamental limits of natural-image-pretrained features on medical imagery.
The resolution thresholds ($\tau_{\text{accept}}$, $\tau_{\text{reject}}$) in the claim resolution rule are fixed; learning these from data could improve adaptability.
Our benchmark evaluates 120 items per domain; larger-scale evaluation would strengthen statistical power.
The autonomous controller currently uses heuristic rules for expert selection; a learned policy could further optimize the cost--accuracy tradeoff.

\paragraph{Future work.}
Extending the expert pool with domain-specific modules (CT windowing, frequency-domain analysis) could address medical domain weaknesses.
Beyond classification, extending the debate to anomaly localization---where the Proposer specifies regions and the Advocate challenges spatial precision---is a natural next step.
Integrating active learning, where uncertain debate outcomes trigger human review and update the system's reference bank, represents a path toward continual improvement in deployed settings.
